\documentclass[11pt]{article}
\usepackage[margin=1in]{geometry}
\usepackage{amsmath,amssymb}
\usepackage{enumitem}
\usepackage{bookmark}
\hypersetup{
  pdftitle={PSET 7: General random variables},
  hidelinks}

\title{PSET 7: General random variables}
\author{}
\date{\vspace{-2.5em}}

% Problems are numbered 1-9 across the document.
\newcounter{problem}
\newcommand{\problem}[1]{\stepcounter{problem}\section*{Problem \theproblem: #1}}

\begin{document}
\maketitle

% Shared assignment questions (header block).
% Include at the top of any problem set with:  % Shared assignment questions (header block).
% Include at the top of any problem set with:  % Shared assignment questions (header block).
% Include at the top of any problem set with:  \input{assignment-qs}
\section*{Assignment Qs}

\begin{itemize}
\item Name
\item How long did this problem set take you (round to nearest hour)? \quad
  0-1 hour \ 2-3 hours \ 4-5 hours \ 5+ hours
\item How difficult was this problem set? \quad
  very easy \ 1 \ 2 \ 3 \ 4 \ 5 \ very challenging
\end{itemize}

\section*{Assignment Qs}

\begin{itemize}
\item Name
\item How long did this problem set take you (round to nearest hour)? \quad
  0-1 hour \ 2-3 hours \ 4-5 hours \ 5+ hours
\item How difficult was this problem set? \quad
  very easy \ 1 \ 2 \ 3 \ 4 \ 5 \ very challenging
\end{itemize}

\section*{Assignment Qs}

\begin{itemize}
\item Name
\item How long did this problem set take you (round to nearest hour)? \quad
  0-1 hour \ 2-3 hours \ 4-5 hours \ 5+ hours
\item How difficult was this problem set? \quad
  very easy \ 1 \ 2 \ 3 \ 4 \ 5 \ very challenging
\end{itemize}


\noindent You will need a calculator (or R) for Problems 3, 4 and 5,
and a standard normal table for Problem 5. Decimal answers rounded to
three places are fine throughout; show the setup (the integral, the
standardization step) in every case.

\problem{Spot the pdf/CDF}

For each function below, state whether it is a valid pdf, a valid CDF,
or neither, and give the specific property that fails when your answer
is ``neither.''

\begin{enumerate}[label=\alph*.]
\item $f(x) = \begin{cases}
  3x^2 & 0 < x < 1 \\
  0 & \text{otherwise}
  \end{cases}$

\item $F(x) = \begin{cases}
  0 & x < 0 \\
  x/6 & 0 \leq x < 6 \\
  1 & x \geq 6
  \end{cases}$

\item $f(x) = \begin{cases}
  x/2 & 0 < x < 3 \\
  0 & \text{otherwise}
  \end{cases}$

\item $F(x) = \begin{cases}
  0 & x < 0 \\
  x^3 & 0 \leq x < 1 \\
  1 & x \geq 1
  \end{cases}$

\item $f(x) = \begin{cases}
  2 - x & 0 < x < 3 \\
  0 & \text{otherwise}
  \end{cases}$

\item For each function you identified as a CDF, give the
  corresponding pdf.

\item Two of the five functions above are a pdf/CDF pair for the same
  random variable. Which two, and how can you tell?
\end{enumerate}

\problem{Identifying the pdf}

A recent college graduate is moving to Des Moines, Iowa to take a new
job, and is looking to purchase a home. When searching for properties
on the real estate websites, it is possible to select the price range
of housing in which one is most interested. Suppose the potential
buyer specifies a price range of \$160,000 to \$340,000, and the
search returns a thousand homes with prices distributed uniformly
throughout that range.\footnote{Inspired by Stinerock 6.1}

\begin{enumerate}[label=\alph*.]
\item Write down $f(x)$, the probability density function associated
  with this random variable.
\item Identify $E[X]$ and $\sigma$.
\end{enumerate}

\problem{Parliamentary elections}

After an election in a parliamentary system, a government (consisting
of a prime minister and a cabinet) is formed by gathering the support
of a majority of newly elected members of parliament. Typically a
government is allowed to remain in power for a certain number of years
before new elections must be called. However, elections can be held
earlier if the Parliament passes a vote of no confidence or the prime
minister decides to dissolve the government. Suppose we are studying
Country Z (which uses a parliamentary system) and we are interested in
the duration of governments. In Country Z, governments must call
elections at least every \textbf{4} years, but they could be called
sooner if there is a vote of no confidence or the prime minister
dissolves the government. Let the continuous random variable $X$
denote the amount of time (measured in years) between the last
election and the calling of the next election. $X$ has support on
all real numbers between 0 and 4. Suppose we know that $X$ has the
probability density function
$$
f(x) = \begin{cases}
kx^4 &  0 < x < 4 \\
0 & \text{otherwise}
\end{cases}
$$
where $k$ is some constant.


\begin{enumerate}[label=\alph*.]
\item Find $k$.
\item Find the CDF of $X$.
\item Find $E(X)$ and $Var(X)$.
\item Find the median of $X$ (the value of $x$ at which
  $\Pr(X \leq x) = \frac{1}{2}$).
\item What is the probability that the government remains in power for
  exactly 4 years? Why?
\item What is the probability that the government remains in power
  between 2 and 4 years? What is the probability that it remains in
  power for less than 2 years or more than 4 years? How are your two
  answers related, and why?
\end{enumerate}

\problem{Calculating the CDF}

$Z$ is distributed according to the following pdf:
$$
f(z) = \begin{cases}
\lambda \exp(-\lambda z) &  0 \le z \\
0 & \text{otherwise}
\end{cases}
$$
(You will sometimes see this written with $\lambda$ replaced by
$1/\beta$.)

\begin{enumerate}[label=\alph*.]
\item What is $F(z)$, the CDF of this distribution? Show the
  integral.
\item Using your answer to the previous question, find
  $\Pr(4 \leq Z \leq 9)$. Leave your answer in terms of $\lambda$.
\item Suppose $\lambda = 2$. Given this, what is $q$, the 25th
  percentile of $Z$?
\item We observe a single random draw from $Z$; what is the
  probability this observation is less than $0.75$? Again suppose
  that $\lambda = 2$.
\end{enumerate}

\problem{Working with normal random variables}

Let $X \sim \text{Normal}(0, 1)$ and
$Y \sim \text{Normal}(2, 9)$.\footnote{Inspired by BT 3.11}

\begin{enumerate}[label=\alph*.]
\item Find $\Pr(X \leq 1.28)$ and $\Pr(X \leq -0.84)$. State the
  property of the normal distribution you used to handle the negative
  value.
\item Find $\Pr(Y \geq 6.5)$, showing the standardization step.
\item Find the distribution of $(Y - 2)/3$, and write down its pdf.
\end{enumerate}

\problem{Relating the Gamma, exponential, and $\chi^2$}

\begin{enumerate}[label=\alph*.]
\item Starting from the Gamma pdf
  $f(x \mid \alpha, \beta) = \frac{\beta^{\alpha}}{\Gamma(\alpha)}
  x^{\alpha - 1} e^{-\beta x}$, show that
  $\text{Gamma}(1, \lambda)$ is the $\text{Exponential}(\lambda)$
  distribution. (Recall $\Gamma(1) = 1$.)
\item Suppose $X \sim \chi^2(6)$. State $E[X]$ and $Var(X)$, and
  explain in one or two sentences where those values come from, using
  the fact that $X = \sum_{i=1}^{n} Z_i^2$ for independent standard
  normal $Z_i$.
\item Suppose $Y_1 \sim \text{Gamma}(2, 3)$ and
  $Y_2 \sim \text{Gamma}(5, 3)$ are independent. What is the
  distribution of $Y_1 + Y_2$? Give its expected value and variance.
\end{enumerate}

\problem{Calculating ideal points}

Suppose An is a voter living in the country of Freedonia, and suppose
that in Freedonia, all sets of public policies can be thought of as
representing points on a single axis (e.g.\ a line running from more
liberal to more conservative).\footnote{Inspired by Grimmer HW11.2} An
has a certain set of public policies that they want to see enacted.
This is represented by point $v$, which we will call An's
\textbf{ideal point}. The utility, or happiness, that An receives from
a set of policies at point $l$ is $U(l) = -(l - v)^2$. In other
words, An is happiest if the policies enacted are the ones at their
ideal point, and they get less and less happy as policies get farther
away from this ideal point. When they vote, An will pick the candidate
whose policies will make them happiest. However, An does not know
exactly what policies each candidate will enact if elected---they have
some guesses, but can't be certain. Each candidate's future policies
can therefore be represented by a continuous random variable $L$ with
expected value $\mu_l$ and variance $Var(L)$.

\begin{enumerate}[label=\alph*.]
\item Express $E(U(L))$ as a function of $\mu_l$, $Var(L)$, and
  $v$. \emph{Hint: expand $(L - v)^2$ and apply the expected value
  rule term by term. You may use $E[L^2] = Var(L) + \mu_l^2$.}

  \vspace{1.5in}

\item Using your answer to (a), explain why we say An is
  \textbf{risk averse}---that is, why An gets less happy as outcomes
  get more uncertain, holding $\mu_l$ fixed.

\item Suppose An is deciding whether to vote for one of two
  candidates: Dwight Schrute or Leslie Knope. An's ideal point is at
  2. Schrute's policies can be represented by a continuous random
  variable $L_S$ with expected value 2 and variance 8, and Knope's
  policies can be represented by a continuous random variable $L_K$
  with expected value 4 and variance 3. Which candidate would An vote
  for, and why? What (perhaps surprising) effect of risk aversion on
  voting behavior does this example demonstrate?
\end{enumerate}

\problem{Moments}

We discussed moments in class and in the slides. Suppose someone were
asking you about our class content. How would you explain the
following in a non-technical way? Approximately 2 sentences each.

\begin{enumerate}[label=\alph*.]
\item What is an expected value, and how does it relate to a moment?
\item What is the first moment, and why do we care about it?
\item What is the second moment, and how does it add additional
  information to the first moment?
\item What do the third and fourth moments describe, and what do they
  tell us that the first two do not?
\end{enumerate}

\problem{Functions}

We discussed the following functions in class: Uniform, Exponential,
Gamma, Normal, $\chi^2$, Student's t.

Provide a one-sentence summary of the general shape and the expected
value of each:

\begin{enumerate}[label=\alph*.]
\item Uniform
\item Exponential
\item Gamma
\item Normal
\item $\chi^2$
\item Student's t
\end{enumerate}

% Shared AI and Resources statement.
% Include in any problem set with:  % Shared AI and Resources statement.
% Include in any problem set with:  % Shared AI and Resources statement.
% Include in any problem set with:  \input{ai-statement}
\section*{AI and Resources statement}

Please list (in detail) all resources you used for this assignment. If
you worked with people, list them here as well. It is not enough to say
that you used a resource for help, you need to be specific on the link
and \emph{how} it was helpful. W/R/T gen AI tools (including GPT,
Claude, etc.) you cannot use them to do work on your behalf---you
cannot put in any of the questions, etc. You can ask for help on logic
/ sample problems. If you do use GPT or other AI tools, you need to
provide a link to your chat transcript. Any suspected academic
integrity violations will be immediately reported to the Dean of
Students.

\section*{AI and Resources statement}

Please list (in detail) all resources you used for this assignment. If
you worked with people, list them here as well. It is not enough to say
that you used a resource for help, you need to be specific on the link
and \emph{how} it was helpful. W/R/T gen AI tools (including GPT,
Claude, etc.) you cannot use them to do work on your behalf---you
cannot put in any of the questions, etc. You can ask for help on logic
/ sample problems. If you do use GPT or other AI tools, you need to
provide a link to your chat transcript. Any suspected academic
integrity violations will be immediately reported to the Dean of
Students.

\section*{AI and Resources statement}

Please list (in detail) all resources you used for this assignment. If
you worked with people, list them here as well. It is not enough to say
that you used a resource for help, you need to be specific on the link
and \emph{how} it was helpful. W/R/T gen AI tools (including GPT,
Claude, etc.) you cannot use them to do work on your behalf---you
cannot put in any of the questions, etc. You can ask for help on logic
/ sample problems. If you do use GPT or other AI tools, you need to
provide a link to your chat transcript. Any suspected academic
integrity violations will be immediately reported to the Dean of
Students.


\end{document}